% Options for packages loaded elsewhere
\PassOptionsToPackage{unicode}{hyperref}
\PassOptionsToPackage{hyphens}{url}
\PassOptionsToPackage{dvipsnames,svgnames,x11names}{xcolor}
\documentclass[
  10pt,
  fontset=fandol]{ctexart}
\usepackage{xcolor}
\usepackage[margin=16mm]{geometry}
\usepackage{amsmath,amssymb}
\setcounter{secnumdepth}{-\maxdimen} % remove section numbering
\usepackage{iftex}
\ifPDFTeX
  \usepackage[T1]{fontenc}
  \usepackage[utf8]{inputenc}
  \usepackage{textcomp} % provide euro and other symbols
\else % if luatex or xetex
  \usepackage{unicode-math} % this also loads fontspec
  \defaultfontfeatures{Scale=MatchLowercase}
  \defaultfontfeatures[\rmfamily]{Ligatures=TeX,Scale=1}
\fi
\usepackage{lmodern}
\ifPDFTeX\else
  % xetex/luatex font selection
\fi
% Use upquote if available, for straight quotes in verbatim environments
\IfFileExists{upquote.sty}{\usepackage{upquote}}{}
\IfFileExists{microtype.sty}{% use microtype if available
  \usepackage[]{microtype}
  \UseMicrotypeSet[protrusion]{basicmath} % disable protrusion for tt fonts
}{}
\makeatletter
\@ifundefined{KOMAClassName}{% if non-KOMA class
  \IfFileExists{parskip.sty}{%
    \usepackage{parskip}
  }{% else
    \setlength{\parindent}{0pt}
    \setlength{\parskip}{6pt plus 2pt minus 1pt}}
}{% if KOMA class
  \KOMAoptions{parskip=half}}
\makeatother
\setlength{\emergencystretch}{3em} % prevent overfull lines
\providecommand{\tightlist}{%
  \setlength{\itemsep}{0pt}\setlength{\parskip}{0pt}}
\usepackage{amsmath,amssymb,mathtools,needspace}
\xeCJKDeclareCharClass{CJK}{"2032,"0400->"04FF,"2160->"216B,"2460->"2473,"25A0,"25C7,"25CB,"25CF}
\IfFontExistsTF{Arial Unicode MS}{\xeCJKsetup{AutoFallBack=true}\setCJKfallbackfamilyfont{\CJKrmdefault}{Arial Unicode MS}}{}
\setcounter{secnumdepth}{0}
\setlength{\emergencystretch}{2em}
\allowdisplaybreaks
\usepackage{bookmark}
\IfFileExists{xurl.sty}{\usepackage{xurl}}{} % add URL line breaks if available
\urlstyle{same}
\hypersetup{
  pdftitle={微积分的逼近法},
  colorlinks=true,
  linkcolor={Maroon},
  filecolor={Maroon},
  citecolor={Blue},
  urlcolor={Blue},
  pdfcreator={LaTeX via pandoc}}

\title{微积分的逼近法}
\author{}
\date{}

\begin{document}
\maketitle

\subsubsection{10.1.1　微积分的逼近法}\label{ux5faeux79efux5206ux7684ux903cux8fd1ux6cd5}

经典数学的基础是微积分。从微积分的观点看，在一切函数中，以多项式最为简单。能否用简单的多项式来逼近一般函数呢？众所周知的 Taylor 分析（1715 年）肯定了这一事实。Taylor 级数

\[
f(x)\sim\sum_{k=0}^{\infty}\frac{f^{(k)}(x_0)}{k!}(x-x_0)^k
\]

表明，一般的光滑函数 \(f(x)\) 可用多项式来近似地刻画。Taylor 分析是 18 世纪初的一项重大的数学成就。

然而 Taylor 分析存在严重的缺陷：它的条件很苛刻，要求 \(f(x)\) 足够光滑并提供出它的各阶导数值 \(f^{(k)}(x_0)\)；此外，Taylor 分析的整体逼近效果差，它仅能保证在展开点 \(x_0\) 的某个邻域内有效。

时移物换，百年之后 Fourier（1822 年）指出，``任何函数，无论怎样复杂，均可表示为三角级数的形式''：

\href{../../source-images/pdf-262.jpeg}{原页}

\[
f(x)\sim\frac{a_0}{2}+\sum_{k=1}^{\infty}(a_k\cos2\pi kx+b_k\sin2\pi kx),\quad 0\leqslant x<1.
\]

这就是今日被称作``Fourier 分析''的数学方法。著名数学家 M. Kline 评价这一数学成就是``19 世纪数学的第一大步，并且是真正极为重要的一步''。

Fourier 关于任意函数都可以表达为三角级数这一思想被誉为``数学史上最大胆，最辉煌的概念''。

Fourier 的成就使人们从 Taylor 分析的理想函数类中解放出来。Fourier 分析不仅放宽了光滑性的限制，还保证了整体的逼近效果。

从数学美的角度来看，Fourier 分析也比 Taylor 分析更美，其基函数系------三角函数系是个完备的正交函数系。尤其值得注意的是，这个函数系可以视作是由一个简单函数 \(\cos x\) 经过简单的伸缩平移变换加工生成的。Fourier 分析表明，任何复杂函数都可以借助于简单函数 \(\cos x\) 来刻画，即

\[
\cos x\xrightarrow{\text{伸缩＋平移}}\text{三角函数系}\xrightarrow{\text{组合}}\text{任意函数 }f(x).
\]

这是一个惊人的事实。在这里，被逼近函数 \(f(x)\) 的``繁''与逼近工具 \(\cos x\) 的``简''两者反差很大，因此 Fourier 逼近很美。Fourier 分析在数学史上被誉``一首数学的诗''，Fourier 则有``数学诗人''的美称。

\end{document}
